\scanpage{17}
- 192 - \hfill n\textsuperscript{o} 391

2) Le réseau $L_1$ est le sous-$\mathrm Z$-module de $L_{\unsure{o|0}}$ formé des $x=\sum_{i=1}^n\xi_i\varepsilon_i$ tels que $\sum_{i=1}^n\xi_i$ soit \uline{pair} ; comme $\xi_i$ et $\xi_i^2$ ont même parité, il revient au même de dire que $(x|x)$ est \uline{pair}. Il est engendré par les $\varepsilon_i\pm\varepsilon_j$, $(i\neq j)$.

3) Le réseau $L_2=L_{\unsure{o|0}}+\mathrm Z.(\frac12\sum_{i=1}^n\varepsilon_i)$ ; il revient au même de dire que les vecteurs $x=\sum_{i=1}^n\xi_i\varepsilon_i$ de $L_2$ sont ceux qui vérifient les relations :
\[
(4)\qquad 2\xi_i\in\mathrm Z\ ,\quad\xi_i-\xi_j\in\mathrm Z\quad\text{quels que soient }i,\ j.
\]
Comme $(\varepsilon_k|\frac12\sum_i\varepsilon_i)=\frac12$ pour tout $k$, et $\|\frac12\sum_i\varepsilon_i\|^2=n/4$, on a $(u|v)\in\frac12\mathrm Z$ pour $u$, $v$ dans $L_2$ si $n$ est pair.

4) Si $n$ est \uline{multiple de} $4$, le réseau $L_3=L_1+\mathrm Z.(\frac12\sum_{i=1}^n\unsure{\xi_i|\varepsilon_i})$, sous-$\mathrm Z$-module de $L_2$; c'est l'ensemble des vecteurs $x=\sum_{i=1}^n\xi_i\varepsilon_i$ de $V$ vérifiant les relations (4) et en outre tels que $\sum_{i=1}^n\xi_i\in2\mathrm Z$ ; on a $(u|v)\in\mathrm Z$ quels que soient $u$, $v$ dans $L_3$.

On notera que le réseau \uline{associé} à $L_{\unsure{o|0}}(V)$ est $L_{\unsure{o|0}}(V^*)$, le réseau \uline{associé} à $L_1(V)$, (resp. $L_2(V)$), est $L_2(V^*)$, (resp. $L_1(V^*)$) ; enfin, le réseau \uline{associé} à $L_3(V)$ est $L_3(V^*)$, car c'est le sous-réseau de $L_2(V^*)$ formé des $x=\sum_{i=1}^n\xi_i\varepsilon_i^*$ tels que $\frac12(\sum_{i=1}^n\xi_i)\in\mathrm Z$, et on vient de voir, (en vertu de l'hypothèse $n\equiv0\ (\operatorname{mod.}4)$) que ces conditions caractérisaient $L_3(V^*)$.

Nous allons voir successivement, dans chacun des 8 n\textsuperscript{os} suivants, comment cette méthode démontre en premier lieu l'existence des systèmes de racines correspondant aux 8 types de schémas de Dynkin autres que $G_2$ ; en outre, pour chaque type de schéma, nous expliciterons :
\mspage{16}
